\section{Related Work}

The literature that precedes this review falls into two strands, and the positioning of this study follows from the gap between them.

\noindent\textbf{Reviews of agents and AI software engineering.} The first strand gathers reviews and visions that map the use of language models and agents in software engineering. Hou et al. carry out a broad systematic review of language-model use across software engineering tasks, organized by lifecycle stage and by technique \cite{hou2024llmse}; He, Treude, and Lo review the internal architecture and cooperation of language-model multi-agent systems, drawing a roadmap toward consolidating them \cite{JundaHe2025}; and Hassan et al. propose a vision of AI-native software engineering, with a challenge roadmap for what they call Software Engineering 3.0 \cite{AhmedEHassan2024}. These contributions are valuable for understanding the anatomy of agents and the direction of the field. But they operate at a level of abstraction that classifies models, cognitive components, and techniques by task, or that projects future scenarios. None characterizes the operational frameworks that, today, organize development work on top of an agent at the level of process and repository.

\noindent\textbf{Process-oriented proposals.} The second strand is made of the frameworks that structure the process themselves. Spec-driven development articulates the specification as a living contract in the age of coding assistants \cite{DeepakBabuPiskala2026}; the distinction between vibe coding and agentic coding formalizes a spectrum of autonomy with practical implications \cite{RanjanSapkota2025}; and frameworks such as MetaGPT, ChatDev, and Reversa instantiate, each in its own way, the organization of the cycle around roles, artifacts, and validation \cite{SiruiHong2023,ChenQian2023,Macedo2026}. Each of these proposals, however, evaluates itself in isolation, with no common reference that would allow them to be compared.

\noindent\textbf{Positioning.} Between a review literature that describes agents abstractly and a proposal literature that evaluates each framework on its own, what is missing is a systematic review that takes the operational frameworks as its object, through a software engineering lens, and that confronts their structure with the available evidence through an explicit reference. That is the contribution of this article: a six-dimension taxonomy validated empirically against a corpus of 37 studies, a thematic synthesis that exposes the field's axes of tension, and a process-oriented research agenda. Unlike practitioner-facing product comparisons, the focus here is the body of literature and the structure that emerges from it, not the recommendation of a specific tool.
